\documentclass[12pt,a4paper]{article}
\usepackage[utf8]{inputenc}
\usepackage{amsmath,amssymb}
\usepackage{geometry}
\usepackage{enumitem}
\usepackage{fancyhdr}
\usepackage{lastpage}

\geometry{a4paper,top=25mm,bottom=25mm,left=25mm,right=25mm}
\setlength{\parindent}{0pt}
\setlength{\parskip}{6pt}

\pagestyle{fancy}
\fancyhf{}
\renewcommand{\headrulewidth}{0pt}
\fancyfoot[L]{Student number:\hspace{4cm}}
\fancyfoot[R]{\thepage\ of \pageref{LastPage}}

\newcommand{\N}{\mathcal{N}}
\newcommand{\KL}{\mathrm{KL}}
\newcommand{\E}{\mathbb{E}}
\newcommand{\ELBO}{\mathrm{ELBO}}

\begin{document}

\begin{center}
{\large\bfseries Technical University of Denmark}\\[4pt]
{\large Advanced machine learning (02460)}\\[4pt]
{\large Practice Exam Set C}\\[4pt]
Exam duration: 2 hours.\quad Aid: None.\\[2pt]
\textbf{IMPORTANT:} Questions are self-contained; all necessary formulas are provided.\\
Only information written on pages 2--11 will be evaluated.\\
Write your student number on all pages at the bottom-left.
\end{center}

\bigskip
\textbf{Contents}
\begin{itemize}[noitemsep]
  \item[\textbf{I}] \textbf{Generative models} \hfill 2
  \begin{itemize}[noitemsep]
    \item Exercise A: Probabilistic PCA \hfill 2
    \item Exercise B: Normalizing flows \hfill 3
    \item Exercise C: Denoising Diffusion Probabilistic Models \hfill 4
  \end{itemize}
  \item[\textbf{II}] \textbf{Representation learning} \hfill 5
  \begin{itemize}[noitemsep]
    \item Exercise D: Pull-back metrics \hfill 5
    \item Exercise E: Geodesic approximation \hfill 6
    \item Exercise F: Curve energy \hfill 7
    \item Exercise G: Neural network reparametrization \hfill 8
  \end{itemize}
  \item[\textbf{III}] \textbf{Graph models} \hfill 9
  \begin{itemize}[noitemsep]
    \item Exercise H: Shallow embeddings \hfill 9
    \item Exercise I: Graph Laplacians and random walks \hfill 10
    \item Exercise J: Graph convolution \hfill 11
  \end{itemize}
\end{itemize}

\newpage
\section*{Part I \quad Generative models}

\subsection*{Exercise A --- Probabilistic PCA}

Consider a probabilistic PCA model with latent variable $z\in\mathbb{R}$ and observed variable
$\mathbf{x}\in\mathbb{R}^2$.  Let $z$ follow a standard Gaussian prior $p(z)=\N(z\mid 0,1)$, and
the likelihood be
\[
  p(\mathbf{x}\mid z) = \N\!\bigl(\mathbf{x}\mid [2,1]^\top z + [0,1]^\top,\; \mathbf{I}_2\bigr).
\]
Recall that if $z\sim\N(0,1)$ and $\mathbf{x}=\mathbf{w}z+\mathbf{b}+\boldsymbol{\varepsilon}$ with
$\boldsymbol{\varepsilon}\sim\N(\mathbf{0},\mathbf{I})$, then
\[
  p(\mathbf{x}) = \N\!\bigl(\mathbf{x}\mid \mathbf{b},\;\mathbf{w}\mathbf{w}^\top+\mathbf{I}\bigr). \tag{1}
\]

\textbf{Question A.1:} What is the marginal distribution $p(\mathbf{x})$?

\begin{itemize}[noitemsep,label={}]
  \item[\textbf{A}] $\square$ \quad $p(\mathbf{x}) = \N\!\bigl(\mathbf{x}\mid [0,1]^\top,\; \begin{bmatrix}5&2\\2&2\end{bmatrix}\bigr)$
  \item[\textbf{B}] $\square$ \quad $p(\mathbf{x}) = \N\!\bigl(\mathbf{x}\mid [0,1]^\top,\; \begin{bmatrix}5&2\\2&3\end{bmatrix}\bigr)$
  \item[\textbf{C}] $\square$ \quad $p(\mathbf{x}) = \N\!\bigl(\mathbf{x}\mid [0,0]^\top,\; \begin{bmatrix}5&2\\2&2\end{bmatrix}\bigr)$
  \item[\textbf{D}] $\square$ \quad $p(\mathbf{x}) = \N\!\bigl(\mathbf{x}\mid [0,0]^\top,\; \begin{bmatrix}5&2\\2&3\end{bmatrix}\bigr)$
\end{itemize}

Give a brief argument for your answer.

\vspace{5cm}

\newpage
\subsection*{Exercise B --- Normalizing flows}

Consider a normalizing flow with $\mathbf{u},\mathbf{x}\in\mathbb{R}^2$, where
$p_{\mathbf{u}}(\mathbf{u})=\N(\mathbf{u}\mid\mathbf{0},\mathbf{I}_2)$ and the transformation
$T:\mathbb{R}^2\to\mathbb{R}^2$ is
\[
  T(\mathbf{u}) = \begin{bmatrix}3&0\\0&2\end{bmatrix}\mathbf{u} + [1,-1]^\top. \tag{2}
\]
The density of $\mathbf{x}=T(\mathbf{u})$ is
\[
  p_{\mathbf{x}}(\mathbf{x}) = p_{\mathbf{u}}(T^{-1}(\mathbf{x}))\cdot|\det J_{T^{-1}}(\mathbf{x})|. \tag{3}
\]

\textbf{Question B.1:} What is $p_{\mathbf{x}}([1,-1]^\top)$?

\begin{itemize}[noitemsep,label={}]
  \item[\textbf{A}] $\square$ \quad $\dfrac{1}{6\pi}$
  \item[\textbf{B}] $\square$ \quad $\dfrac{1}{12\pi}$
  \item[\textbf{C}] $\square$ \quad $\dfrac{1}{18\pi}$
  \item[\textbf{D}] $\square$ \quad $\dfrac{1}{36\pi}$
\end{itemize}

Give a brief argument for your answer.

\vspace{6cm}

\newpage
\subsection*{Exercise C --- Denoising Diffusion Probabilistic Models}

Consider a DDPM with $T=5$ steps.  The variance schedule is
\[
  \beta_1=\tfrac{1}{4},\quad\beta_2=\tfrac{1}{3},\quad\beta_3=\tfrac{1}{2},\quad
  \beta_4=\tfrac{1}{4},\quad\beta_5=\tfrac{1}{4}.
\]
Recall the forward process in eq.~(3), the marginal distribution
\[
  q(\mathbf{x}_t\mid\mathbf{x}_0) = \N(\mathbf{x}_t\mid\sqrt{\bar\alpha_t}\,\mathbf{x}_0,\,(1-\bar\alpha_t)\mathbf{I}),
  \quad \bar\alpha_t=\prod_{\tau=1}^t(1-\beta_\tau), \tag{4}
\]
and the \emph{forward posterior}
\[
  q(\mathbf{x}_{t-1}\mid\mathbf{x}_t,\mathbf{x}_0) = \N\!\bigl(\mathbf{x}_{t-1}\mid\tilde\mu_t,\;\tilde\beta_t\mathbf{I}\bigr), \tag{5}
\]
where
\[
  \tilde\mu_t = \frac{\sqrt{\bar\alpha_{t-1}}\,\beta_t}{1-\bar\alpha_t}\,\mathbf{x}_0
              + \frac{\sqrt{1-\beta_t}\,(1-\bar\alpha_{t-1})}{1-\bar\alpha_t}\,\mathbf{x}_t,
  \qquad
  \tilde\beta_t = \frac{\beta_t(1-\bar\alpha_{t-1})}{1-\bar\alpha_t}. \tag{6}
\]

\textbf{Question C.1:} Assuming $\mathbf{x}_0=[0,0]^\top$ and $\mathbf{x}_2=[3,0]^\top$,
what is $q(\mathbf{x}_1\mid\mathbf{x}_2,\mathbf{x}_0)$?

First compute $\bar\alpha_1$ and $\bar\alpha_2$.  Then use equations (5) and (6).

\vspace{2cm}

\[
  \bar\alpha_1 = \hspace{4cm} \bar\alpha_2 =
\]

\vspace{2cm}

\[
  q(\mathbf{x}_1\mid\mathbf{x}_2=[3,0]^\top,\;\mathbf{x}_0=[0,0]^\top) =
\]

Sketch your derivations below.

\vspace{5cm}

\newpage
\section*{Part II \quad Representation learning}

\subsection*{Exercise D --- Pull-back metrics}

Consider $f:\mathbb{R}^2\to\mathbb{R}^3$ defined by
\[
  f(\mathbf{x}) = \bigl[x_1^2+x_2,\;\; x_1+x_2^2,\;\; x_1 x_2\bigr]^\top. \tag{7}
\]

\textbf{Question D.1:} Compute the Jacobian $J_{\mathbf{x}}$ of $f$ at $\mathbf{x}$, and then
derive the pull-back metric $G_{\mathbf{x}} = J_{\mathbf{x}}^\top J_{\mathbf{x}}\in\mathbb{R}^{2\times 2}$.

\vspace{7cm}

\textbf{Question D.2:} Evaluate $G_{\mathbf{x}}$ at $\mathbf{x}=[1,0]^\top$.

\[
  G_{[1,0]^\top} =
  \begin{bmatrix} & \\ & \end{bmatrix}
\]

\newpage
\subsection*{Exercise E --- Geodesic approximation}

Consider a generative model $\mathbf{y}=f(\mathbf{x})$, where $\mathbf{x}\in\mathbb{R}^d$ and
$\mathbf{y}\in\mathbb{R}^D$ with $d<D$, and $f:\mathbb{R}^d\to\mathbb{R}^D$ is smooth.  We wish to
find the \emph{geodesic} (shortest path on the manifold) between $\mathbf{c}_0$ and
$\mathbf{c}_{N+1}$ by minimising over intermediate nodes $\mathbf{c}_{1:N}$:
\[
  \mathbf{c}^*_{1:N} = \operatornamewithlimits{argmin}_{\mathbf{c}_{1:N}}
    \sum_{n=1}^{N+1} \|f(\mathbf{c}_n)-f(\mathbf{c}_{n-1})\|. \tag{8}
\]

\textbf{Question E.1:} Is the piecewise-linear latent path with nodes
$\{\mathbf{c}_0,\ldots,\mathbf{c}_{N+1}\}$ a discrete approximation of the shortest
path in the \emph{latent space} $\mathbb{R}^d$, or the shortest path on the
\emph{manifold} $f(\mathbb{R}^d)\subset\mathbb{R}^D$?

\begin{itemize}[noitemsep,label={}]
  \item[\textbf{A}] $\square$ \quad It approximates the shortest path in the latent space $\mathbb{R}^d$.
  \item[\textbf{B}] $\square$ \quad It approximates the shortest path on the manifold $f(\mathbb{R}^d)$.
\end{itemize}

Explain your reasoning.

\vspace{3cm}

\textbf{Question E.2:} A student proposes to instead minimise
$\sum_{n=1}^{N+1}\|{\mathbf{c}_n-\mathbf{c}_{n-1}}\|$ (straight lines in latent space).
Would this produce the same result as eq.~(8) in general?

\begin{itemize}[noitemsep,label={}]
  \item[\textbf{A}] $\square$ \quad Yes, they are always equivalent.
  \item[\textbf{B}] $\square$ \quad No, they differ in general.
\end{itemize}

Give a brief argument.

\vspace{3cm}

\newpage
\subsection*{Exercise F --- Curve energy}

Consider the generative model $\mathbf{y}=f(\mathbf{x})$, where
\[
  f(\mathbf{x}) = \bigl[x_1,\;\; x_1^2 + x_2^2,\;\; x_2\bigr]^\top. \tag{9}
\]
Consider the two latent curves $\mathbf{c}_1(t)=t[1,0]^\top$ and $\mathbf{c}_2(t)=t[0,1]^\top$
for $t\in[0,1]$.

\textbf{Question F.1:} Derive the pull-back metric $G_{\mathbf{x}}=J_{\mathbf{x}}^\top J_{\mathbf{x}}$.

\vspace{5cm}

\textbf{Question F.2:} Evaluate $E(\mathbf{c}_1)$ and $E(\mathbf{c}_2)$ using
$E(\mathbf{c})=\int_0^1\dot{\mathbf{c}}^\top G_{\mathbf{c}(t)}\dot{\mathbf{c}}\;\mathrm{d}t$.

\[
  E(\mathbf{c}_1) = \hspace{6cm} E(\mathbf{c}_2) =
\]

Sketch your derivations below.

\vspace{4cm}

\newpage
\subsection*{Exercise G --- Neural network reparametrization}

Consider the neural network $f:\mathbb{R}\to\mathbb{R}$
\[
  f(x\mid w_1, w_2) = w_1\,\sigma(w_2\, x), \tag{10}
\]
where $\sigma(x)=\max(0,x)$ is the ReLU activation function.

\textbf{Question G.1:} Let $w_3 = \alpha w_1$ for some $\alpha>0$.  Derive a value for $w_4$ such
that
\[
  f(x\mid w_1, w_2) = f(x\mid \alpha w_1, w_4) \quad \text{for all } x. \tag{11}
\]

\[
  w_4 =
\]

Sketch your derivations below.

\vspace{4cm}

\textbf{Question G.2:} Discuss the implication of your answer for the statistical identifiability
of the parameter vector $(w_1, w_2)$: can we uniquely determine $(w_1,w_2)$ from observations
$\{(x_i, f(x_i))\}$?

\vspace{4cm}

\textbf{Question G.3:} Would the same issue arise if $\sigma$ were a \emph{bijective} activation
function (e.g., $\sigma(x)=x$)?  Give a brief argument.

\vspace{3cm}

\newpage
\section*{Part III \quad Graph models}

\subsection*{Exercise H --- Shallow embeddings}

We use the \emph{sigmoid decoder}
$P(A_{uv}=1) = \sigma(\mathbf{z}_u^\top\mathbf{z}_v + b)$ where $\sigma(x)=1/(1+e^{-x})$, in
a 2-dimensional embedding space with $\mathbf{z}_v=[x,y]^\top$.

\textbf{Question H.1:} Let $\mathbf{z}_u=[1,0]^\top$ and $b=0$.  Where could node $v$ be
embedded so that $P(A_{uv}=1)=0.5$?

\begin{itemize}[noitemsep,label={}]
  \item[\textbf{A}] $\square$ \quad On the line $x=0$ (the $y$-axis)
  \item[\textbf{B}] $\square$ \quad On the line $x=1$
  \item[\textbf{C}] $\square$ \quad On a circle of radius 1 centred at the origin
  \item[\textbf{D}] $\square$ \quad On the line $y=0$
  \item[\textbf{E}] $\square$ \quad None of the above
\end{itemize}

Show the main steps of your derivation.

\vspace{4cm}

\textbf{Question H.2:} Two nodes $u$ and $v$ are embedded as $\mathbf{z}_u=[2,0]^\top$ and
$\mathbf{z}_v=[0,2]^\top$.  With $b=0$, which node pair has higher probability of an edge
according to the sigmoid decoder?

\begin{itemize}[noitemsep,label={}]
  \item[\textbf{A}] $\square$ \quad Node pair $(u, u)$, i.e.\ a self-loop on $u$
  \item[\textbf{B}] $\square$ \quad Node pair $(u, v)$
  \item[\textbf{C}] $\square$ \quad Both pairs have equal probability
  \item[\textbf{D}] $\square$ \quad Cannot be determined without knowing $b$
\end{itemize}

Give a brief argument.

\vspace{3cm}

\newpage
\subsection*{Exercise I --- Graph Laplacians and random walks}

Consider a graph whose random-walk Laplacian $L_{\mathrm{rw}}=I-D^{-1}A$ is
\[
  L_{\mathrm{rw}} =
  \begin{bmatrix}
    1 & -\tfrac{1}{3} & -\tfrac{1}{3} & -\tfrac{1}{3} \\[4pt]
    -1 & 1 & 0 & 0 \\[4pt]
    -1 & 0 & 1 & 0 \\[4pt]
    -1 & 0 & 0 & 1
  \end{bmatrix},
\]
with nodes labelled $a, b, c, d$ in row order.

\textbf{Question I.1:} Sketch the graph.

\vspace{3cm}

\textbf{Question I.2:} A random walk starts at node $b$.  What is the probability $p$ of being
at node $a$ after exactly 2 steps?

\[
  p =
\]

Give a brief argument.

\vspace{3cm}

\textbf{Question I.3:} The stationary distribution $\boldsymbol{\pi}$ of the random walk
satisfies $D^{-1}A\,\boldsymbol{\pi}=\boldsymbol{\pi}$ (equivalently, $L_{\mathrm{rw}}\boldsymbol{\pi}=\mathbf{0}$).
Without full computation, write down the stationary distribution $\boldsymbol{\pi}$ up to
normalisation.  Recall that for an undirected graph, $\pi_v\propto d_v$.

\[
  \boldsymbol{\pi} \propto
\]

\vspace{2cm}

\newpage
\subsection*{Exercise J --- Graph convolution}

Consider the same 4-node star graph from Exercise I (centre $a$, leaves $b, c, d$).
Its adjacency matrix and eigendecomposition are:
\[
  A = \begin{bmatrix}0&1&1&1\\1&0&0&0\\1&0&0&0\\1&0&0&0\end{bmatrix},
  \qquad
  \lambda_1=-\sqrt{3},\;\lambda_2=\lambda_3=0,\;\lambda_4=\sqrt{3}.
\]
A graph convolution is $\mathbf{y}=\sum_{k=0}^N h[k]\,A^k\,\mathbf{x}$.

\textbf{Question J.1:} Using the filter $h[0]=1, h[1]=-1$ (and $h[k]=0$ otherwise) and the
signal $\mathbf{x}=[0,1,1,1]^\top$, compute the filtered signal $\mathbf{y}$.

\[
  \mathbf{y} = [\phantom{00},\phantom{00},\phantom{00},\phantom{00}]^\top
\]

Show the main steps.  (\emph{Hint:} Compute $A\mathbf{x}$ directly from the adjacency structure.)

\vspace{5cm}

\textbf{Question J.2:} The Fourier-domain filter matrix for filter $h[k]$ is
$M=\sum_{k=0}^N h[k]\Lambda^k$.  For $h[0]=1, h[1]=-1$, what is the first diagonal entry
$M_{11}$ (corresponding to $\lambda_1=-\sqrt{3}$)?

\[
  M_{11} =
\]

\end{document}
